% !TeX root = ../../金融经济学习题讲义.tex
% 源 PDF：references/金融经济学/史上最全最完整的金融经济学答案(王江版).pdf，页 14–18（印刷页 13–17）

\zhchapter{套利和资产定价}

第 3 章在 Arrow--Debreu 市场结构下讨论均衡与效率；本章回到第 2 章的一般支付矩阵 $X$，讨论\textbf{证券价格的基本性质}、\textbf{定价原理}、\textbf{无套利}与\textbf{资产定价基本定理}。

\begin{tcolorbox}[knowledge,title=4.1 一般市场结构]
多数可交易证券在多个状态下都有正支付，称为\textbf{复合证券}（相对只在单一状态付 1 的阿罗证券而言）。设 1 期有 $\Omega$ 个状态，市场上有 $N$ 只证券（$n=1,\dots,N$），第 $n$ 只证券的支付向量为 $X_n=\begin{bmatrix}x_{1,n} \\ \vdots \\ x_{\Omega,n}\end{bmatrix}$，\textbf{市场结构矩阵}为（与第 2.3 节一致，行为状态、列为证券）：
\[
    X=[X_{\cdot,1},\,\cdots,\,X_{\cdot,N}]=
    \begin{bmatrix}
        x_{1,1}   & \cdots & x_{1,N}   \\
        \vdots    & \ddots & \vdots    \\
        x_{\omega,1} & \cdots & x_{\omega,N} \\
        \vdots    & \ddots & \vdots    \\
        x_{\Omega,1} & \cdots & x_{\Omega,N}
    \end{bmatrix}.
\]

\tcbline

\textbf{A.\ 冗余证券}\quad 若证券 $j$ 的支付列 $x_j$ 可由其余证券的\textbf{复制组合}（replicating portfolio）线性表出，即存在 $\theta_{\setminus j}$ 使得
\[
    x_j=X_{\setminus j}\theta_{\setminus j},
\]
则称 $j$ 为\textbf{冗余证券}。冗余证券不增加可达的 1 期支付集合；在无摩擦时，市场结构 $X$ 与删去第 $j$ 列后的 $X_{\setminus j}$ 在资源配置意义上\textbf{等价}，均衡与分配分析中常可忽略冗余列。\textcolor{red}{\textbf{更精确的说法：}}存在或去掉某证券\textbf{不改变均衡配置}时，该证券才是冗余的。存在\textbf{市场摩擦}时，用其他证券复制 $x_j$ 可能额外花费，表观冗余的证券在交易中仍可能不可替代。

\tcbline

\textbf{B.\ 证券市场的等价描述}\quad 无摩擦时，可用\textbf{线性无关}的 $N$ 只证券描述市场，$\mathrm{rank}(X)=N$（通常设 $N\le\Omega$）。任取 $N$ 个线性无关的组合向量 $\theta_1,\dots,\theta_N$，令 $H=[\theta_1,\cdots,\theta_N]$（$N\times N$ 可逆），则新支付矩阵 $X_\theta=XH$ 与原市场\textbf{等价}：原证券列空间不变，$X=X_\theta H^{-1}$，每只原证券都可由新组合复制。

\tcbline

\textbf{C.\ 生成（span）与完全市场}\quad 记 $\mathbf{1}_\omega$ 为第 $\omega$ 分量为 1、其余为 0 的列向量（状态 $\omega$ 阿罗证券的支付）。若 $\mathrm{rank}(X)=\Omega$ 且 $N\ge\Omega$，则对每一 $\omega$ 存在组合 $\theta_\omega$ 使 $X\theta_\omega=\mathbf{1}_\omega$。当 $X$ 为 $\Omega\times\Omega$ 可逆方阵时，解唯一：
\[
    \theta_\omega=X^{-1}\mathbf{1}_\omega.
\]
此时称不确定性由证券\textbf{生成}（span）；市场与第 3 章完全阿罗市场等价。

\begin{tcolorbox}[knowledge]
\textbf{定理 4.1}\quad 无摩擦时，证券市场\textbf{完全}当且仅当线性无关证券只数等于状态数（可取 $N=\Omega$ 且 $\mathrm{rank}(X)=\Omega$，此时 $X$ 可逆）。这与定义 3.2、\textcolor{red}{\textbf{行满秩}}判据一致。
\end{tcolorbox}
\end{tcolorbox}

\begin{tcolorbox}[knowledge,title=4.2 套利]
记交易证券的 0 期价格向量为 $S=\begin{bmatrix}S_1 \\ \vdots \\ S_N\end{bmatrix}$。组合 $\theta=\begin{bmatrix}\theta_1 \\ \vdots \\ \theta_N\end{bmatrix}$ 的 0 期成本为 $S^{\mathsf T}\theta$，1 期支付为 $X\theta$。从支付 $X$ 到价格 $S$ 的映射称为\textbf{资产定价关系}（资产定价模型）。若 $X\theta\ge 0$，称该组合在 1 期满足\textbf{有限责任}。

\begin{tcolorbox}[knowledge]
\textbf{定义 4.1}\quad 组合 $\theta$ 构成\textbf{套利机会}，若同时满足：
\begin{enumerate}
    \item $S^{\mathsf T}\theta\le 0$（0 期净支出非正）；
    \item $X\theta\ge 0$（1 期各状态支付非负）；
    \item 上述两式中\textbf{至少有一个严格不等式}成立（0 期严格收入或某状态严格正支付）。
\end{enumerate}
等价地，扩展支付向量 $\begin{bmatrix}-S^{\mathsf T}\theta \\ X\theta\end{bmatrix}$ 非负且至少一分量严格为正。
\end{tcolorbox}

常分为三类（记严格不等方向与教材一致）：
\begin{itemize}
    \item \textbf{第 1 类：}$S^{\mathsf T}\theta<0$ 且 $X\theta=0$——0 期白捡钱、未来无负债；
    \item \textbf{第 2 类：}$S^{\mathsf T}\theta=0$ 且 $X\theta>0$——\textbf{套利组合}（零成本、未来严格正支付）；
    \item \textbf{第 3 类：}$S^{\mathsf T}\theta<0$ 且 $X\theta>0$——0 期收入为正且未来仍为正。
\end{itemize}

\begin{tcolorbox}[example]
\textbf{例 4.1}\quad $S^{\mathsf T}=[1,\,1,\,2]$，
$X=\begin{bmatrix}1&0&2\\1&2&0\\1&2&0\end{bmatrix}$（三证券在各状态下的支付）。组合 $\theta_1=[2;\,-1;\,-1]$：$S^{\mathsf T}\theta_1=-1$，$X\theta_1=0$，为\textbf{第 1 类套利}。经济直觉：支付 $[2;2;2]$ 可由 2 单位证券 1 以成本 $2$ 复制，也可由 1 单位 $B$ 加 1 单位 $C$ 以成本 $3$ 复制——同支付不同价。
\end{tcolorbox}
\end{tcolorbox}

\begin{tcolorbox}[knowledge,title=4.3 无套利原理]
套利机会只依赖公开的支付 $X$ 与价格 $S$，不依赖私人信息或对状态概率的主观信念。

\begin{tcolorbox}[knowledge]
\textbf{定理 4.2}\quad 市场均衡中\textbf{不存在}套利机会。
\end{tcolorbox}

\textbf{直观：}若存在套利 $\theta$，在\textbf{不满足性}下参与者可无限放大该组合，得到严格更优的消费，与均衡最优矛盾。\textbf{前提：}至少一名参与者不满足；市场\textbf{无摩擦}（无成本、无限制、价格接受者等，见第 2.3 节）。均衡 $\Rightarrow$ 无套利是\textbf{必要}条件；无套利 $\not\Rightarrow$ 均衡（还需出清与偏好等）。

\begin{tcolorbox}[knowledge]
\textbf{定义 4.2}\quad 证券市场上不存在套利机会，称为满足\textbf{无套利原理}（no-arbitrage principle）。
\end{tcolorbox}
\end{tcolorbox}

\begin{tcolorbox}[knowledge,title=4.4 资产定价基本定理（开端）]
将证券（或组合）的 1 期支付 $x\in\mathrm{R}^{\Omega}$ 映射为 0 期价格 $S$，记\textbf{定价算子}（估值算子）$V$：
\[
    S=V(x).
    \tag{4.1}
\]
对列向量 $x$ 为某只证券支付时，$S_n=V(X_n)$。

\begin{tcolorbox}[knowledge]
\textbf{定理 4.3（一价定律）}\quad 若两只证券（或组合）未来支付相同，$x_1=x_2$，则 0 期价格必相同：$V(x_1)=V(x_2)$。
\end{tcolorbox}

无套利蕴含一价定律；反之一般不成立。进一步可证：\textbf{无套利}等价于存在\textbf{状态价格}（或正线性定价泛函）使 $S=X^{\mathsf T}\phi$（完全市场时 $\phi$ 唯一，不完全市场时 $\phi$ 可能不唯一甚至某些分量可为负，见练习 4.5）。鞅测度表述见练习 4.8。
\end{tcolorbox}


\begin{problem}
\problemno{4.1} 经济在 1 期有 4 个可能状态。在市场中有 5 只可交易证券，它们的支付矩阵 $X$ 如下：
\[
    X=\begin{bmatrix}
        1 & 1 & 0 & 0 & 0 \\
        1 & 2 & 1 & 0 & 0 \\
        1 & 3 & 2 & 1 & 0 \\
        1 & 4 & 3 & 2 & 1
    \end{bmatrix}.
\]
\begin{enumerate}[label=(\alph*)]
    \item 证明市场是完全的。
    \item 证明状态价格向量存在且唯一。并给出它的值。
    \item 证明市场上存在冗余证券。
    \item 选择一组足以保证市场完全的复合证券。把每一只 Arrow-Debreu 证券都表示成这些证券的组合。称这些组合为状态或有组合。
    \item 用状态或有组合把任一由五只交易证券组成的组合表示出来。
    \item 证明状态或有组合可生成状态空间。
\end{enumerate}
\end{problem}
\begin{solution}
记 5 只交易证券的支付向量分别为
\[
    X_{1}=\begin{bmatrix}1 \\ 1 \\ 1 \\ 1\end{bmatrix},\quad X_{2}=\begin{bmatrix}1 \\ 2 \\ 3 \\ 4\end{bmatrix},\quad X_{3}=\begin{bmatrix}0 \\ 1 \\ 2 \\ 3\end{bmatrix},\quad X_{4}=\begin{bmatrix}0 \\ 0 \\ 1 \\ 2\end{bmatrix},\quad X_{5}=\begin{bmatrix}0 \\ 0 \\ 0 \\ 1\end{bmatrix}
\]
\begin{enumerate}[label=(\alph*)]
    \item 因为 $X_{1}$、$X_{2}$、$X_{3}$、$X_{4}$ 的支付是互相独立的四只证券，而状态数等于 4，因而证券市场是完全的；
    \item 假设状态价格向量为 $\phi=\begin{bmatrix}\phi_{1} \\ \phi_{2} \\ \phi_{3} \\ \phi_{4}\end{bmatrix}$，那么有：
    \[
        \begin{aligned}
            \phi_{1}+\phi_{2}+\phi_{3}+\phi_{4}&=1\\
            \phi_{1}+2\phi_{2}+3\phi_{3}+4\phi_{4}&=2.5\\
            \phi_{2}+2\phi_{3}+3\phi_{4}&=1.5\\
            \phi_{3}+2\phi_{4}&=0.75\\
            \phi_{4}&=0.25
        \end{aligned}
    \]
    求解上述方程组可以得到唯一的解为 $\phi=\begin{bmatrix}\tfrac14 \\ \tfrac14 \\ \tfrac14 \\ \tfrac14\end{bmatrix}$；
    \item 买入 1 单位的证券 2 卖出 1 单位的证券 3，得到的支付是 $\begin{bmatrix}1 \\ 1 \\ 1 \\ 1\end{bmatrix}$，恰好就是证券 1 的支付，因而市场上存在冗余证券；
    \item 证券 2、3、4、5 可以保证市场完全，而且我们有
    \[
        \begin{aligned}
            X_{2}-2X_{3}+X_{4}&=\begin{bmatrix}1 \\ 0 \\ 0 \\ 0\end{bmatrix}\\
            X_{3}-2X_{4}+X_{5}&=\begin{bmatrix}0 \\ 1 \\ 0 \\ 0\end{bmatrix}\\
            X_{4}-2X_{5}&=\begin{bmatrix}0 \\ 0 \\ 1 \\ 0\end{bmatrix}\\
            X_{5}&=\begin{bmatrix}0 \\ 0 \\ 0 \\ 1\end{bmatrix}
        \end{aligned}
    \]
    \item 假设组合中 5 只交易证券的头寸分别为 $\theta_{1}$、$\theta_{2}$、$\theta_{3}$、$\theta_{4}$、$\theta_{5}$，那么状态或有组合的头寸分别为 $\theta_{1}+\theta_{2}$，$\theta_{1}+2\theta_{2}+\theta_{3}$，$\theta_{1}+3\theta_{2}+2\theta_{3}+\theta_{4}$，$\theta_{1}+4\theta_{2}+3\theta_{3}+2\theta_{4}+\theta_{5}$；
    \item 因为状态或有组合的支付是 $\mathrm{R}^{4}$ 空间的一组基，因而状态或有组合可生成支付空间。
\end{enumerate}
\end{solution}

\begin{problem}
\problemno{4.2} 银行 1 的存贷利率为 $r_{1}$ 而银行 2 的为 $r_{2}$，且 $r_{1}>r_{2}$。证明存在套利机会。
\end{problem}
\begin{solution}
从银行 2 借钱并全部存入银行 1，参与者就可以获得第 2 类套利。
\end{solution}

\begin{problem}
\problemno{4.3} 银行 $i$ 的贷款利率为 $r_{ia}$（即顾客借钱的利率）而存款利率为 $r_{ib}$，$i=1,2$。
\begin{enumerate}[label=(\alph*)]
    \item 证明无套利要求 $r_{ia}\ge r_{ib}$。
    \item 两个银行的存贷利率应该满足什么样的条件才能使市场上不存在套利机会。
\end{enumerate}
\end{problem}
\begin{solution}
\begin{enumerate}[label=(\alph*)]
    \item 反证法。如果 $r_{ia}<r_{ib}$，那么从银行 $i$ 借入 1 块钱并存入银行 $i$，那么在 0 期支付为 0 而 1 期支付为 $r_{ib}r_{ia}>0$，参与者可获得第 2 类套利。因此，$r_{ia}\ge r_{ib}$；
    \item $r_{1a}>r_{2b}$，$r_{1b}<r_{2a}$。
\end{enumerate}
\end{solution}

\begin{problem}
\problemno{4.4} 考虑练习 4.1 中定义的市场。假设 5 只交易证券的价格分别为 $1$，$2$，$3/2$，$1$，$1/2$。
\begin{enumerate}[label=(\alph*)]
    \item 在这样的价格和支付下存在套利机会吗？
    \item 你能构建第 1、2 以及 3 类套利机会吗？
\end{enumerate}
\end{problem}
\begin{solution}
\begin{enumerate}[label=(\alph*)]
    \item 买入 1 单位的证券 4，卖出 2 单位的证券 5，0 期支付为 0，而 1 期支付为 $\begin{bmatrix}0 \\ 0 \\ 1 \\ 0\end{bmatrix}>0$，因而存在套利机会。
    \item 上面构造了第 2 类套利机会。卖出 1 单位的证券 1 和 1 单位的证券 3，同时买入 1 单位的证券 2，那么 0 期支付为 $\tfrac12$，而 1 期支付为 0，这是第 1 类套利机会。同时进行上述两种交易就可以得到第 3 类套利机会。
\end{enumerate}
\end{solution}

\begin{problem}
\problemno{4.5} 在 1 期有三个可能状态：$\{\omega_{1},\omega_{2},\omega_{3}\}$。在经济中有三只交易证券，$a$，$b$，$c$，它们的支付向量如下：
\[
    X_{a}=\begin{bmatrix}1\\1\\0\end{bmatrix},\quad
    X_{b}=\begin{bmatrix}0\\1\\1\end{bmatrix},\quad
    X_{c}=\begin{bmatrix}1/2\\1\\1/2\end{bmatrix}.
\]
这些证券的 0 期均衡价格为：$S_{c}=S_{b}=S_{c}=1$。令 $\phi=\begin{bmatrix}\phi_{1} \\ \phi_{2} \\ \phi_{3}\end{bmatrix}$ 为状态价格向量。
\begin{enumerate}[label=(\alph*)]
    \item 证明 $\begin{bmatrix}2 \\ -1 \\ 2\end{bmatrix}$ 是一个可能的状态价格向量，因为它给出了三支现有证券的价格。
    \item “资产定价基本定理”指出状态价格必须是非负的。在一个两期经济里，对于证券市场这意味着什么？并给出其解释。
    \item 上面的例子是这个原理的反例吗？为什么？
    \item 证明 $\begin{bmatrix}1/2 \\ 1/2 \\ 1/2\end{bmatrix}$ 是另外一个可能的状态价格向量。在这种情况下，所有状态价格都是正的。
    \item 解释为什么存在多个可能的状态价格向量，它们对已有证券都给出相同的价格。
    \item 考虑一个标的资产为证券 $a$、执行价格为 $1/2$ 的欧式买权。这份期权的支付是多少？你能为它定价吗？
\end{enumerate}
\end{problem}
\begin{solution}
\begin{enumerate}[label=(\alph*)]
    \item 因为
    \[
        2\times 1+(-1)\times 1=1,\quad (-1)\times 1+2\times 1=1,\quad 2\times\tfrac12+(-1)\times 1+2\times\tfrac12=1
    \]
    因而 $\begin{bmatrix}2 \\ -\tfrac12 \\ 2\end{bmatrix}$ 是一个可能的状态价格向量；
    \item 状态价格向量的非负性意味着不存在套利机会；
    \item 不是。因而它只是可能的状态价格之一，并不一定是真实的状态价格；
    \item 因为
    \[
        \tfrac12\times 1+\tfrac12\times 1=1,\quad \tfrac12\times 1+\tfrac12\times 1=1,\quad \tfrac12\times\tfrac12+\tfrac12\times 1+\tfrac12\times\tfrac12=1
    \]
    因而 $\begin{bmatrix}\tfrac12 \\ \tfrac12 \\ \tfrac12\end{bmatrix}$ 是一个可能的状态价格向量；
    \item 因为市场是不完全的，因而可能存在多个状态价格，它们都给出已有证券的价格。
\end{enumerate}
\end{solution}

\begin{problem}
\problemno{4.6} 远期合约（Forward）是在未来（即 1 期）以当期（即 0 期）确定的价格交易某一证券或商品的合约，而在当期没有现金流产生。在当期约定的未来交易价格称作远期价格。比方说，考虑一份在 1 期以 $F_{0}$ 的价格购买 1 盎司黄金的远期合约。合约的买方同意在 1 期（合约的到期日）从卖方买入 1 盎司黄金。在当期没有现金交换。令 $S_{t}$ 为 $t=0,1$ 期黄金的价格，而 $r_{F}$ 为利率。考虑在 1 期得到 1 盎司黄金的两种方法：
\begin{itemize}
    \item 现在购买一份远期合约。在 1 期，支付远期价格 $F_{0}$ 并获得黄金。
    \item 现在支付 $S_{0}$ 买入 1 盎司黄金，并且持有它直到 1 期。假设从现在起持有黄金直到 1 期没有另外的成本或收益。
\end{itemize}
在这两种情况下，我们在 1 期同样得到 1 盎司黄金，而那时它的价值为 $S_{1}$。
\begin{enumerate}[label=(\alph*)]
    \item 由于这两种策略给予我们相同的支付 $S_{1}$，无套利原理意味着它们的成本必须一样。用当期的价格 $S_{1}$ 和 $r_{F}$ 给出远期价格。
    \item 假设持有 1 盎司黄金不仅得到它的 1 期出售收益，在 1 期还得到额外收益 $y$。$y$ 将如何影响远期价格？
\end{enumerate}
\end{problem}
\begin{solution}
\begin{enumerate}[label=(\alph*)]
    \item 第一种方法的成本是 $\dfrac{F_{0}}{1+r_{F}}$，而第二种方法的成本是 $S_{0}$，根据无套利原理，$\dfrac{F_{0}}{1+r_{F}}=S_{0}$，因此，$F_{0}=(1+r_{F})S_{0}$；
    \item 由于第二种方法还得到额外收益，由无套利原理原理有 $F_{0}<(1+r_{F})S_{0}$。（如果收益是确定性的，那么应该有 $F_{0}+y=(1+r_{F})S_{0}$）
\end{enumerate}
\end{solution}

\begin{problem}
\problemno{4.7} 英镑现时对美元的价格（即英镑/美元汇率）为 $S_{0}$（即一英镑的美元价格为本 $S_{0}$），美元利率为 $r_{A}$ 而英镑利率为 $r_{B}$。考虑一份在 1 期购买 1 英镑的远期合约。远期价格（即远期汇率）应该是多少？提示：考虑在 1 期得到 1 英镑的两种策略。
\begin{itemize}
    \item 购买一份远期合约。在 1 期，支付远期价格 $F_{0}$ 购入 1 英镑。
    \item 现在以当前汇率 $S_{0}$ 买入 $x$ 英镑并以 $r_{B}$ 的利率进行投资，在 1 期得到 1 英镑。
\end{itemize}
得到的远期价格、现时汇率和利率之间的关系就叫作抛补利率平价理论（Covered Interest Rate Parity）。
\end{problem}
\begin{solution}
第一种策略的 0 期支付是 $\dfrac{F_{0}}{1+r_{A}}$ 美元，而第二种策略的 0 期支付是 $\dfrac{S_{0}}{1+r_{B}}$ 美元，根据无套利原理，$\dfrac{F_{0}}{1+r_{A}}=\dfrac{S_{0}}{1+r_{B}}$，即 $F_{0}=\dfrac{1+r_{A}}{1+r_{B}}S_{0}$。
\end{solution}

\begin{problem}
\problemno{4.8} 考虑练习 4.1 定义的市场。假设证券 1，2，5 的价格是 1，2，$1/2$。
\begin{enumerate}[label=(\alph*)]
    \item 计算状态价格和利率。
    \item 构建等价鞅测度。
    \item 证明 5 只证券的价格都具有鞅性质。
\end{enumerate}
\end{problem}
\begin{solution}
\begin{enumerate}[label=(\alph*)]
    \item 假设状态价格向量为 $\phi=\begin{bmatrix}\phi_{1} \\ \phi_{2} \\ \phi_{3} \\ \phi_{4}\end{bmatrix}$，那么有：
    \[
        \begin{aligned}
            \phi_{1}+\phi_{2}+\phi_{3}+\phi_{4}&=0.8\\
            \phi_{1}+2\phi_{2}+3\phi_{3}+4\phi_{4}&=2\\
            \phi_{2}+2\phi_{3}+3\phi_{4}&=1.2\\
            \phi_{4}&=0.2
        \end{aligned}
    \]
    求解上述方程组可以得到唯一的解为 $\phi=\begin{bmatrix}0.2 \\ 0.2 \\ 0.2 \\ 0.2\end{bmatrix}$。因此，证券 4 的价格为 $0.6$，而利率为 $r_{F}=\dfrac{1}{\phi_{1}+\phi_{2}+\phi_{3}+\phi_{4}}-1=0.25$；
    \item 等价鞅测度为 $Q=\begin{bmatrix}0.25 \\ 0.25 \\ 0.25 \\ 0.25\end{bmatrix}$；
    \item 对于证券 1，$\dfrac{\mathrm{E}^{Q}[X_{1}]}{1+r_{F}}=\dfrac{1}{1+0.25}=0.8$，正好就是它的价格；类似可以证明其它证券的价格也有鞅性质。
\end{enumerate}

Notes：证券价格不完全，现在改为证券 1，2，3，5 的价格分别为 0.8，2，1.2，0.2。
\end{solution}
